\documentclass{article}
\PassOptionsToPackage{numbers,compress}{natbib}
\usepackage[preprint]{neurips_2025}
\usepackage[utf8]{inputenc}
\usepackage[T1]{fontenc}
\usepackage{hyperref}
\usepackage{url}
\usepackage{booktabs}
\usepackage{amsmath,amsfonts}
\usepackage{microtype}
\usepackage{graphicx}
\usepackage{xcolor}
\newcommand{\TODO}[1]{\textcolor{red}{[TODO: #1]}}

\title{Generating Paintings with a Spectrally-Normalized DCGAN on WikiArt}

\author{%
  \TODO{Your Name} \\
  ECE 285, UC San Diego \\
  \texttt{\TODO{email}}
}

\begin{document}
\maketitle

\begin{abstract}
We train a deep convolutional GAN (DCGAN) with a spectrally-normalized discriminator on all 81{,}444 paintings of the WikiArt dataset at $64\times64$ resolution. We report the hyperparameters, training dynamics, and sample quality, measured with FID and Inception Score computed using PyTorch-Ignite. \TODO{one-sentence headline result: final FID / IS}
\end{abstract}

\section{Method}
\textbf{Data.} The Kaggle WikiArt dataset contains 81{,}444 paintings across 27 styles. We center-crop each image to a square, resize it to $64\times64$ with Lanczos filtering, and scale pixels to $[-1,1]$. During training we apply random horizontal flips. Style labels are not used, so the GAN is unconditional.

\textbf{Architecture.} The generator follows DCGAN~\cite{radford2016dcgan}. A 128-d Gaussian latent passes through a $4\times4$ transposed convolution, then three stride-2 transposed-conv/BatchNorm/ReLU blocks (512$\to$256$\to$128$\to$64 channels), then a final transposed convolution with $\tanh$ (3.8M parameters). The discriminator mirrors it with four stride-2 convolutions (64$\to$512 channels), LeakyReLU(0.2), and a $4\times4$ convolution that outputs one logit (2.8M parameters). Every discriminator layer uses spectral normalization~\cite{miyato2018sn} in place of BatchNorm.

\textbf{Objective.} We use the hinge loss: $L_D=\mathbb{E}[\max(0,1-D(x))]+\mathbb{E}[\max(0,1+D(G(z)))]$ and $L_G=-\mathbb{E}[D(G(z))]$. Samples and metrics use an exponential moving average (EMA) of the generator weights.

\begin{table}[h]
\centering
\caption{Hyperparameters.}
\label{tab:hparams}
\small
\begin{tabular}{ll@{\hspace{2em}}ll}
\toprule
Resolution & $64\times64$ & Optimizer & Adam ($\beta_1{=}0.5$, $\beta_2{=}0.999$) \\
Batch size & 128 & Learning rate (G / D) & $2\times10^{-4}$ / $2\times10^{-4}$ \\
Latent dim & 128 & D steps per G step & 1 \\
Epochs & 100 (63.6k iterations) & Generator EMA decay & 0.999 \\
Loss & hinge & Precision & bf16 autocast \\
\bottomrule
\end{tabular}
\end{table}

\textbf{Evaluation.} We follow the PyTorch-Ignite GAN evaluation guide and use \texttt{ignite.metrics.FID} and \texttt{InceptionScore} inside an Ignite \texttt{Engine}. Generated and real images are mapped to $[0,1]$, bilinearly upsampled to $299\times299$, and ImageNet-normalized for torchvision's InceptionV3. FID is computed on 2048-d pool features between 10{,}000 generated images and 10{,}000 random training images. \TODO{training time on RTX 5070 laptop GPU}

\section{Results}
\TODO{loss curve figure + discussion}

\TODO{FID/IS table (and FID vs epoch figure)}

\TODO{sample grid figure}

\TODO{discussion: what works, failure modes}

{\small
\begin{thebibliography}{9}
\bibitem{radford2016dcgan} A.~Radford, L.~Metz, and S.~Chintala. Unsupervised representation learning with deep convolutional generative adversarial networks. \emph{ICLR}, 2016.
\bibitem{miyato2018sn} T.~Miyato, T.~Kataoka, M.~Koyama, and Y.~Yoshida. Spectral normalization for generative adversarial networks. \emph{ICLR}, 2018.
\end{thebibliography}
}
\end{document}
